
\ChapterWithSubtitle{Oceans and Ice}{something}
%\chapter{Heating by Radiation}

%%%%%%%%%%%%%%%%%%%%%%%%%%%%%%%%%%%%%%%%%%%%%
\section*{Equipment List}

\begin{multicols}{2}   % change 2 -> 3 for three columns
\begin{equipment}
\item Ruler
\item Stopwatch
\item Digital Thermometer
\item 250 mL beaker
\item 500 mL beaker
\item crushed ice
\item large ice cube
\item hot plate
\item aquarium
\item test tube (20mL)
\item test tube stopper and capillary tube
\end{equipment}
\end{multicols}


%%%%%%%%%%%%%%%%%%%%%%%%%%%%%%%%%%%%%%%%%%%%%
\section{Introduction}

Today we will consider two physical principles connected to the effects of climate change on sea level: (i) floating, and (ii) thermal expansion of fluids.  On the first topic, we will try to understand the difference between floating ice (icebergs) and land-based ice, and their effects on the oceans when melting.  On the second topic, we will measure the small effect of water's expansion when heated, and estimate it's effect on sea-level rise.

\subsection{Buoyant Force}

Early in the semester we introduce Newton's Second Law:
	%
	\begin{equation*}
	\sum \vec{F} = m \vec{a} \quad \quad 
	\left(
	\begin{array} {c}
	 \text{The sum of the (vector) forces on an object of} \\
	\text{ mass $m$ cause changes to its (vector) velocity }.
	\end{array}
	\right)
	\end{equation*}
	%
where, recall that the acceleration, $a = \frac{\Delta \vec{v}}{\Delta t}$, expresses the change in velocity of the object. We will apply this law today only to objects in {\bf equilbrium}, where $\vec{a} = 0$, meaning that the sum of the forces in each situation will add up to zero.  

\noindent To understand floating, we will need a new type of force: the {\bf bouyant force}.  The origin of this force is quite easy to understand.  Recall that we talked about how the pressure within a fluid increases with depth below the surface --- this is felt as an effect on your ear drums when you go to the bottom of the swimming pool. You are probably also familiar with the fact that if you descend into a pool of water, you will generally be pushed back up to the surface.  This is due to the so-called bouyant force, which arises as the sum of all the pressure forces acting on your body.  We can derive a simple expression for this force by imagining, instead of your body, a ``you-shaped parcel'' of water submerged in the pool.  Because the outline of that you-shaped parcel is the same as you, it will experience exactly the same pressure forces that you had felt when submerged.  But while you were pushed to the surface, the net effect on the parcel of water is to hold it in place, in equilibrium.  Specifically, the weight of the water in the parcel (the force of gravity on it), must be exactly equal to the sum of those pressure forces.  \\[-6pt]

\noindent In the space below, draw two pictures side by side: you submerged in the pool of water and a you-shaped parcel of water (indicated with a dotted-line outline?) submerged in a pool of water.  Draw little vector arrows pushing in on the surface of each object to represent the pressure forces.  

\answerspace[5cm]

\noindent Looking at your picture, and understanding what we do about pressure in a fluid, why do you think that the net effect of the pressure forces is to push the object up?

\answerspace[5cm]

\noindent From this analysis, we can immediately write down an equation for the {\bf strength of the bouyant force in water}:
	%
	\begin{equation*}
	F_B = \text{weight of water displaced by object} = m_\text{disp} \, g
	\end{equation*}
	%
where, in the second equality, $m_\text{disp}$ is the mass of the water displaced and we have used the expression for weight: $W = F_g = m g$, with $g = \unit{10}{\meter/\second^2}$.\\[-6pt]

\noindent Considering this equation, why is your submerged body pushed to the surface of the water (particularly if you have a lungful of air). In other words, why is the bouyant force up on you larger than the gravity force down on you?

\answerspace

\subsection{Floating and icebergs}

Let's now consider the case of an object that is not submerged, but is floating at the surface of the water.  For it to float, there must be some part of the object that is submerged (to get displaced water and a bouyant force), but there is also some part of it that is above the surface of the water.  Draw a picture of a box floating in water. Let it have a cross-sectional surface area  $A$ (the top of the box) and a height of $h$.  The height of the submerged part of the box is $d$.  Draw two force vectors on the box: force of gravity down and bouyant force up.  

\answerspace

\noindent Newton's Second Law for this box says that the force up must cancel the force down (because it is in equilibrium,  not accelerating).  Therefore you can write an equation that says ``the force of gravity, $F_g$, is equal to the bouyant force, $F_B$''.  Do this below. Then, in the next line down, replace each side of the equation with the expression for that force (both have been given above).

\answerspace

\noindent Now we will introduce the {\bf density} to replace masses with volumes.  The definition of density, $\rho$, of an object is:
	%
	\begin{equation*}
	\rho = \frac{m_\text{obj}}{V_\text{obj}} = \frac{\text{mass of the object}}{\text{volume of the object}}
	\end{equation*}
	%
where the Greek letter $\rho$ (``rho'') is conventionally used as the symbol for density.  Use this to re-write your Newton's Second Law equation above.  You should first solve the density definition for mass and then use that expression on both sides of your equation.

\answerspace

\noindent We will come together as a class to discuss this equation.  Then use the euqation to answer two questions: (1) Which objects cannot float?  (2) How submerged is  {\em floating ice}, which has a density of $\rho_\text{ice} = \frac{9}{10} \,\rho_\text{water}$?

\answerspace

%%%%%%%%%%%%%%%%%%%%%%%%%%%%%%%%%%%%%%%%%%%%%
\subsection{Thermal expansion of water}

Most materials expand when heated. This includes almost all solids\footnote{One important exception is water.  We've seen above that ice is {\em less dense} than liquid water.  Above 4$\degree$C, liquid water does expand with temperature (and become less dense), but below 4$\degree$C water {\em expands with cooling}. This was likely important for the evolution of life, since pools of water freeze from the top down, leaving liquid water below.}, and almost all liquids and gases.  An increased ocean temperature due to climate change (primarily affecting the surface, the top \unit{100}{\meter}) will cause that water to expand in volume.  This, together with melting ice on land, will be the cause of global sea level rise over the next centuries.\\[-6pt]

\noindent 

While solid materials can be understood to expand in length when heated --- a ``bimetallic strip'' is a commonly used element of a thermostat to convert temperature change into mechanical movement --- the only meaningful type of expansion for liquids and gases is {\em volumetric} expansion.  This is generally assumed to be a proportionality between the fractional volume increase with the temperature (over a small range of temperatures):
	%
	\begin{equation*}
	\frac{\Delta V}{V} = \beta \, \Delta T
	\end{equation*}
	%
where the quantity, $\beta$ (Greek letter ``beta''), is the {\em coefficient of expansion}.  Inspecting this equation, what units does $\beta$ have?

\answerspace[5cm]

\noindent A table of thermal expansion coefficients for water is listed below\footnote{For sources of these numbers see, e.g., \href{https://web.archive.org/web/20250207061706/https://engineeringtoolbox.com/water-density-specific-weight-d_595.html}{this website}, Table A3.1 in this \href{https://www.aos.wisc.edu/~aos610/Gill_App3.pdf}{book chapter}, and \href{https://www.science.org/doi/10.1126/sciadv.abq0793}{this article}.}.  In today's lab we will attempt to measure the thermal expansion of water in the range of 5--50$\degree$C. Notice that there is not a single value here, but we should expect to see a range as temperature is varied.

\vspace{1cm}

\begin{center}
\begin{tabular}{|l|c|}
\hline
substance (temperature) & $\beta$ ($10^{-6}/\degree$C)\\
\hline
\hline
water (0.1$\degree$C) & -68\\
\hline
seawater, salinity 3.5\% (0.1$\degree$C) & +30\\
\hline
water (1$\degree$C) & -50\\
\hline
water (4$\degree$C) & 3\\
\hline
seawater (4$\degree$C) & 102\\
\hline
water (10$\degree$C) & 88\\
\hline
seawater (10$\degree$C) & 167\\
\hline
water (15$\degree$C) & 151\\
\hline
water (20$\degree$C) & 207\\
\hline
water (25$\degree$C) & 257\\
\hline
seawater (25$\degree$C) & 297\\
\hline
water (30$\degree$C) & 303\\
\hline
seawater (30$\degree$C) & 341\\
\hline
water (50$\degree$C) & 454\\
\hline
\end{tabular}
\end{center}
%
%\begin{tabular}{|l|c|}
%\hline
%T ($\degree$C) & $\beta$ ($10^{-6}/\degree$C)\\
%\hline
%\hline
%water (0.1$\degree$C) & -68\\
%\hline
%water (1$\degree$C) & -50\\
%\hline
%water (4$\degree$C) & 3\\
%\hline
%seawater (4$\degree$C) & 102\\
%\hline
%water (10$\degree$C) & 88\\
%\hline
%seawater (10$\degree$C) & 167\\
%\hline
%water (15$\degree$C) & 151\\
%\hline
%water (20$\degree$C) & 207\\
%\hline
%water (25$\degree$C) & 257\\
%\hline
%seawater (25$\degree$C) & 297\\
%\hline
%water (30$\degree$C) & 303\\
%\hline
%seawater (30$\degree$C) & 341\\
%\hline
%water (50$\degree$C) & 454\\
%\hline
%\end{tabular}%
%\hfill %
%\begin{tabular}{|l|c|}
%\hline
%T ($\degree$C)  & $\beta$ ($10^{-6}/\degree$C)\\
%\hline
%\hline
%water (0.1$\degree$C) & -68\\
%\hline
%seawater, salinity 3.5\% (0.1$\degree$C) & +30\\
%\hline
%water (1$\degree$C) & -50\\
%\hline
%water (4$\degree$C) & 3\\
%\hline
%seawater (4$\degree$C) & 102\\
%\hline
%water (10$\degree$C) & 88\\
%\hline
%seawater (10$\degree$C) & 167\\
%\hline
%water (15$\degree$C) & 151\\
%\hline
%water (20$\degree$C) & 207\\
%\hline
%water (25$\degree$C) & 257\\
%\hline
%seawater (25$\degree$C) & 297\\
%\hline
%water (30$\degree$C) & 303\\
%\hline
%seawater (30$\degree$C) & 341\\
%\hline
%water (50$\degree$C) & 454\\
%\hline
%\end{tabular}

\newpage


%%%%%%%%%%%%%%%%%%%%%%%%%%%%%%%%%%%%%%%%%%%%%
\section{Experiments and Analysis}

{\bf \underline{Safety Notes}:} 
	\begin{itemize}
	\item As always, take care using the hotplate: be aware of when it is hot and keep your neighbors aware.  
	\item For the thermal expansion experiment, we will be using some {\bf very fragile glass capillary tubes}.  Handle these with extreme care as they could break and stab you!
	\end{itemize}

\subsection{Floating and Melting Ice}

Here we will study the effects on melting ice on ``sea'' level (the water level in your aquarium).  

\noindent {\bf Procedure:} Use the blank space below to make diagrams and to record measurements and observations.
	%
	\begin{itemize}
	\item There are small ``icebergs'' in the cooler for you to use, but you should leave them there until they are needed. They are approximately cubes of length \unit{5}{\centi\meter}.
	\item Place the can in your aquarium --- this will serve as your land mass ---  and then fill the aquarium with water to an appropriate level.  It should be high enough that the ice cube can float, but does not need to be much higher than that.  When filling with water, use {\em cold water}, so: fill your 500~mL beaker with (crushed) ice first, then top it off with tap water and mix.  Record the temperature of your water.
	\item Measure the size of your aquarium. Then remove the can briefly, measure the height of the water, and calculate the volume of the water.  Replace the can.
	\item With a whiteboard marker, mark the initial water level and record that initial water volume and the height value.
	\item Obtain your iceberg and quickly measure it to get its volume.
	\item Place the iceberg in the water. Record your observations. Mark your new water level and record its height.
	\item Place the iceberg on the can (land). Record your observations.  Mark the water level and record the height.
	\item Remove the iceberg from the aquarium. Place it in a dry 500~mL beaker and heat it on the hotplate.
	\end{itemize}
	%
	
\vspace{10cm}

\noindent {\bf Analysis and Summary questions}: Record your responses and calculations below.
	%
	\begin{enumerate}
	\item How does the presence of iceberg in the water change the water level?  Is the volume difference of the ``sea'' more or less than the volume of the iceberg?
	\item What will happen to the water level when an ``ice sheet'' on ``land'' melts?  By how will the volume of the sea increase, how much will its level rise?
	\item Which has a bigger effect on sea level, dropping an ice mass into the water, or melting an ice mass that was initially on land.
	\item Pour your melted icecube into the aquarium to check your conclusions.
	\item Why, theoretically, should we expect the outcome you found? [Hint: Consider the mass of the iceberg and the mass of the water that it creates when melted.  These must be the same, by the principle of conservation of matter.  Write down an equation $m_\text{ice} = m_\text{water,melt}$, and then use the density equation to solve for the volume of melted water.  You result should look familiar.]
	\item Given the current state of Earth --- with icebergs and a floating (northern) polar ice cap, and ice on the (continental) southern pole --- which ice, when melted, will cause more sea rise?
	\end{enumerate}
	
	
%%%%%%%%%%%%%%%%%%%%%%%%%%%%%%%%%%%%%%%%%%%%%
\newpage
\subsection{Thermal Expansion of Water}

By heating water that is trapped in a test tube and connected narrow glass ``capillary tube'', we will attempt to observe the expansion of water when heated.

\noindent {\bf Procedure:} Use the blank space below to make diagrams and to record measurements and observations. Open a Google Sheets document to record your temperature and height measurements.
	%
	\begin{itemize}
	\item Our heating setup will be the test tube, with stopper and capillary tube, inside of the 250~mL beaker (acting as a heat bath), on top of the hot place.  We would like to record values for a wide range of temperatures, so you should first prepare very cold water in your 250~mL beaker using crushed ice. Preheat the hotplate at 150$\degree$C.
	\item Fill the test tube with that same cold water (pouring out of the beaker into the test).  
	\item Now you want to place the capillary tube into the water and get some water to flow up into the tube.  You should be able to do this by pushing the stopper into the test tube. The height of the water in the tube will be controlled by how far you push the stopper in.  You might need to experiment to get: (i) the water in the tube above the stopper, (ii) the water level not far above the stopper,  (iii) no extraneous bubbles or water in the tube, (iv) the water level is steady.
	\item Record the initial height of your water column in the capillary tube.
	\item Place the test tube in the water-filled 250~mL beaker and place that on the hotplate.  
	\item Take measurements of temperature and height of water column every minute or so, until the temperature has reached $\sim50\degree$C.
	\end{itemize}

\vspace{10cm}



\noindent {\bf Analysis and Summary questions}: Record your responses and calculations below.
	%
	\begin{itemize}
	\item Using your Google Sheet, make a new column next to ``height'' for the volume of water in the capillary tube (the tube is circular with an inner diameter of \unit{1}{\milli\meter}).
	\item Make a quick plot of volume vs.\ temperature to see where you have good data.   There might be some scatter to the initial point, so you want to avoid that. 
	\item Make another column for $\Delta V$, which should be zero at your initial height and increasing as height increases.
	\item Make a plot (either in Google Sheets or by hand) of $\Delta V$ vs $\Delta T$.  If you made it in Google Sheets, print out the plot.
	\item By hand: select a region of values where you see approximately linear data. Fit only those points to a straight line.
	\item Calculate the slope of your line.
	\item Using the equation provided in the introduction, calculate the ``expansion coefficient'', $\beta$.  You will need to know that your test tube has a volume of \unit{20}{\milli\liter}.
	\end{itemize}



 